% ============================================================================
\section{Implementation and Supporting Derivations}
% ============================================================================

\subsection{Configuration}
\label{app:reference_tables}

\begin{table}[H]
\centering
\footnotesize
\caption{Configuration of the four-domain study. Entries marked TODO
require the corresponding training or evaluation artifact.}
\label{tab:pipeline}
\begin{tabular}{p{0.19\linewidth}p{0.73\linewidth}}
\toprule
Component & Configuration \\
\midrule
Student & Qwen3-1.7B \citep{qwen2025qwen3}, non-thinking mode;
\missing{frozen model and tokenizer revisions} \\
Teachers & Domain-specific Qwen3-1.7B RL checkpoints for math and
instruction following; one shared Qwen3-4B weight set for code and science.
All use the same vocabulary and tokenizer.
\missing{exact checkpoint identifiers, RL training recipes, and scores} \\
Training data & $16{,}000$ unique prompts per task. Joint runs consume
shuffled streams without replacement; single-task references reshuffle
only after exhausting the first cycle.
\missing{datasets, versions, licenses, filters, and held-out split} \\
Default loss & Teacher top-$64$ corrected reverse KL
(Equation~\ref{eq:dense_topk_loss}); full-vocabulary probabilities;
position mean within response, response mean within task; $w_i=1/4$ \\
Rollout & One response per prompt and one update per fresh batch;
maximum response length $4{,}096$; truncated responses retained.
\missing{temperature, EOS handling, and decoding settings} \\
Step and budget & Joint runs: $b=4$, $G=16$, $2\leq m_i\leq8$,
$500$ updates and $32{,}000$ attempted responses per run. Uniform and
per-step GPAS: three training seeds; seven other configurations: one
exploratory seed each. Four single-task references: one seed each,
$G=16$ on the active task and $500$ updates \\
Optimizer & AdamW;
\missing{learning rate and schedule, betas, epsilon, weight decay,
clipping, precision, and memory configuration} \\
GPAS & Pre-step bias-corrected Adam second moment for $D$;
noise and time EMA decay $0.9$; first step Uniform;
exact feasible-integer enumeration for both per-step and cost-aware
counts \\
Checkpoint diagnostics & All three Uniform seeds at updates $50$, $250$,
and $450$; seed-1 update $50$ sets the prediction rule. Calibration:
$32$ independent micro-batches per domain split into two halves for
cross-fitted margins; independent trial and evaluation data.
Precise gradient: $32$ micro-batches per domain; equal-count sweep:
$m_i\in\{2,4,8,16,32\}$ outside the main run's count caps \\
Evaluation & $128$ held-out prompts per task every $50$ steps;
MATH-500, LiveCodeBench, IFBench, and GPQA-Diamond.
\missing{versions, evaluation seeds, and decoding protocols} \\
System & One 96\,GB learner GPU and one 48\,GB rollout/teacher GPU;
the code and science endpoints share weights.
\missing{GPU models, inference and training software revisions,
peak memory, and GPU hours} \\
\bottomrule
\end{tabular}
\end{table}

\subsection{Dense loss implementation and online statistics}
\label{app:implementation}

At each student-generated prefix, the teacher returns the identifiers
and full-vocabulary log probabilities of its top-$64$ tokens. The
student computes its full-vocabulary log normalizer and gathers its
probabilities at those identifiers. Autodifferentiation is applied to
the whole expression in Equation~\ref{eq:dense_topk_loss}, including
the $-p_\theta$ term. The retained probabilities are not renormalized
within the top-$k$ set, and the loss does not use a sampled tail
correction. The operation is dense over the retained set and uses
sparse teacher output.

GPAS collects statistics from unweighted micro-batch gradients. For a
task block, initialize its running mean $h=0$ and scalar $S=0$.
After micro-batch $s$, form $\delta=g_{i,s}-h$, then update
\begin{equation}
    h\leftarrow h+\delta/s,\qquad
    S\leftarrow S+\frac{s-1}{s}\|D\delta\|_2^2.
    \label{eq:welford}
\end{equation}
At the end of the block, $\widehat e_i=S/(m_i-1)$ and
$A\leftarrow A+w_i h$. The mean buffer is reused for the next task.
The same recurrence with $D=I$ records unpreconditioned noise.
Reductions use float32 accumulation. Gradients are unscaled before
measurement if mixed-precision loss scaling is active, and clipping
is applied to the assembled gradient $A$.

The optimizer state is fixed during this collection. Before an Adam
second moment is available, the first Uniform step uses $D=I$.
Subsequent steps use the pre-step bias-corrected second moment.
The first observation initializes each noise EMA; later observations
use decay $0.9$. The scheduler reads the completed step's averages
before generating the next step's data.

For both allocation modes, enumerate
$\mathcal M=\{m\in\mathbb Z^M: \sum_i m_i=G,\,
m_{\min}\leq m_i\leq m_{\max}\}$ and evaluate the objective using
the previous step's smoothed statistics. Per-step GPAS minimizes
$\sum_i w_i^2\widetilde e_i/m_i$; the cost-aware mode minimizes its
product with the estimated step time. Ties first minimize distance
to Uniform and then follow lexicographic task order. If all noise
estimates are zero, use Uniform. This finite search is exact for the
estimated objective and stated integer constraints. Largest-remainder
rounding of the continuous solution does not in general solve that
integer problem. No current-step observations are used to select
the counts of the same step.

\begin{figure}[t]
    \centering
    \includegraphics[width=\linewidth]{figures/gpas_main_figure.pdf}
    \caption{GPAS combines task means with fixed weights $w$ while
    allocating counts $m$ from observed gradient variation and, in its
    cost-aware mode, measured time. The pre-step scaling $D$ defines the
    variance statistic. The profiles are schematic and do not report
    measured MOPD results or certify joint task progress.}
    \label{fig:gpas_overview}
\end{figure}

\subsection{Comparison recipes}
\label{app:baselines}

\paragraph{D$^3$-MOPD scheduling.}
We use the remaining-gap and descent-velocity construction of
\citet{sun2026d3mopd}. Let $\bar L_i(t)$ be the smoothed task loss
and $L_i^0$ the mean of its first five observations. With $W=10$,
$R'\leq R=3$ the number of available windows, and
$\epsilon_{\rm KL}=0.15$, the composite signal is
\begin{equation}
\begin{split}
    r_i(t)&=\bar L_i(t)/\max\{L_i^0,\epsilon_{\rm KL}\},\\
    \nu_i(t)&=\max\left\{0,\frac1{R'}
       \sum_{j=0}^{R'-1}
       \frac{\bar L_i(t-(j+1)W)-\bar L_i(t-jW)}
            {\max\{\bar L_i(t-(j+1)W),\epsilon_{\rm KL}\}}\right\},\\
    s_i(t)&=r_i(t)\nu_i(t).
\end{split}
\end{equation}
We retain the published max-normalization, softmax temperature $0.5$,
per-domain probability floor $0.10$, and multiplicative batch jitter
of amplitude $0.30$. The scheduler refreshes every $10$ steps,
using the paper's history warmup and EMA window of $10$.
All KL denominators use the published floor of $0.15$.
The signal is computed from the default dense task loss in our pipeline.
The resulting mixture is converted to micro-batch counts with the
shared bounds $[2,8]$. Thus this is a controlled adaptation to our loss
scale, micro-batch unit, and budget.

The two D$^3$ rows use the same scheduling rule. The fixed-weight
row assembles $A=\sum_iw_i\bar g_i$; the recipe-level row assembles
$A=\sum_i(m_i/G)\bar g_i$. This makes the effect of task weighting
visible alongside the comparison of scheduling signals.

\paragraph{Open-MOPD with one update per rollout.}
We follow the released student top-$16$ dense candidate loss,
token-share balancing, and gap-aware weighting
\citep{gao2026openmopddiagnosingfixingcapability,openmopd2026code}.
The released implementation multiplies candidate-wise detached
advantages by candidate-wise student probability ratios and then
sums over candidates. Each retained token therefore contributes its own gradient. In our setting, the initial task target is uniform, and each
batch contains four micro-batches per task. The recipe's token and
gap weights remain active. One optimizer step per rollout batch
corresponds to $K=1$, so there is no reuse-induced reward staleness.
\missing{Record the exact configuration for gap smoothing, floors,
and weight normalization from the integrated implementation.}

\paragraph{Tail-aware supervision.}
TA-OPD \citep{huang2026tailaware} groups the omitted vocabulary
into one tail event. For teacher support $S_i(c)$, write
$p_{\rm tail}=1-\sum_{a\in S_i(c)}p_\theta(a\mid c)$ and
$q_{\rm tail}=1-\sum_{a\in S_i(c)}q_i(a\mid c)$. Its position loss is
\begin{equation}
    \ell_i^{\rm TA}=
       \sum_{a\in S_i(c)}p_\theta(a\mid c)
                     \log\frac{p_\theta(a\mid c)}{q_i(a\mid c)}
       +p_{\rm tail}\log\frac{p_{\rm tail}}{q_{\rm tail}}.
\end{equation}
Both runs use $k=64$, the same response averaging and fixed task
weights, and the same training budget as the core comparison.
Only the allocation differs between TA-OPD + Uniform and
TA-OPD + GPAS.

\subsection{Frozen objectives, sampling, and state-distribution change}
\label{app:local_objective}

The task loss $L_i^t$ in Equation~\ref{eq:local_task_loss} holds the
student's step-$t$ rollout distribution $\mathcal R_i^t$ fixed while
varying $\theta$. Under differentiation--expectation interchange,
$\mu_i=\nabla L_i^t(\theta_t)$ and
$\mu_w=\nabla F_t(\theta_t)$. Counts chosen from history before sampling
preserve this weighted mean when the conditional task sampling
law is unchanged. This is a sequence of frozen-distribution objectives;
it does not differentiate through the policy that generated the prefixes.

Equation~\ref{eq:policy_drift} separates optimization on the old states
from the change in visited states. More explicitly, if
$h_i(\theta;x,y)$ is the response-mean retained-support loss, its second
term is
\begin{equation}
    \int h_i(\theta_{t+1};x,y)\,
          d(\mathcal R_i^{t+1}-\mathcal R_i^t)(x,y).
\end{equation}
This term has no general sign. For example, lower moving-policy loss
could reflect a shift toward prefixes on which agreement was already
higher. Fixed-reference loss, fresh-policy loss, and downstream
capability are therefore separate measurements. A finite-$k$ loss
measures the stated retained-support surrogate, not unbiased
full-vocabulary KL.

\paragraph{Finite prompt streams.}
The independent micro-batch covariance model is exact for the independent
diagnostic draws, conditional on the fixed state. Joint training consumes
prompts without replacement, so its batch means are correlated. To state
the correction, consider a uniform sample of $n_i=bm_i$ distinct prompts
from a remaining task pool of size $N_i$. At a fixed student state, let
$u_i(x)$ be the response-averaged expected gradient for prompt $x$,
$R_i(x)$ the conditional rollout-gradient covariance, $S_{u,i}$ the
finite-population covariance of $u_i$ with denominator $N_i-1$, and
$\bar R_i$ the population average of $R_i(x)$. With independent response
randomness conditional on prompts,
\begin{equation}
    \operatorname{Cov}(\bar g_i)
      =\frac1{n_i}\left(1-\frac{n_i}{N_i}\right)S_{u,i}
        +\frac1{n_i}\bar R_i.
    \label{eq:finite_population_covariance}
\end{equation}
The first term follows from finite-population sampling; the second
follows by conditioning on the selected prompts and averaging their
independent rollout noise. For independent task pools, the covariance of
$A$ is the sum of these matrices weighted by $w_i^2$. The local progress
and Cantelli arguments continue to hold with this covariance, but the
simple $\Sigma_i/m_i$ formula is then an approximation. GPAS uses that
approximation during training. Logs must record remaining pool sizes so
its accuracy can be checked; the finite-population correction is small
only when the sampled fraction and the relevant prompt-mean contribution
make it small. Exhausting and reshuffling a specialist's stream also
changes the conditional prompt pool and must be recorded.

\subsection{Mean progress and directional reliability}
\label{app:joint_progress}

All quantities below are conditional on the same checkpoint, chosen
counts, and fixed positive diagonal $D$. Suppose the $g_{j,s}$ are
independent with mean $\mu_j$ and covariance $\Sigma_j$. Then
\begin{equation}
    \mathbb E A=\mu_w,\qquad
    \Sigma(m):=\operatorname{Cov}(A)
      =\sum_j\frac{w_j^2}{m_j}\Sigma_j.
    \label{eq:combined_covariance}
\end{equation}
Linearity gives $\mathbb E X_i=\mu_i^\top D\mu_w=a_i$, and the
variance of a linear functional gives
$\operatorname{Var}(X_i)=\mu_i^\top D\Sigma(m)D\mu_i=v_i(m)$.
The first-order loss change is $-\eta X_i$. Consequently, if
$a_i\leq0$, its expectation is nonnegative for every admissible count
vector. Additional samples can change the variance and finite-step
curvature contribution, but cannot make this expected first-order
change negative at the fixed mean and $D$.

For completeness, write $Z_i=a_i-X_i$, so that
$\mathbb E Z_i=0$ and $\mathbb E Z_i^2=v_i$. For any $\lambda\geq0$
and $a_i>0$, Markov's inequality implies
\begin{equation}
    \Pr(X_i\leq0)=\Pr(Z_i\geq a_i)
      \leq\frac{\mathbb E(Z_i+\lambda)^2}{(a_i+\lambda)^2}
      =\frac{v_i+\lambda^2}{(a_i+\lambda)^2}.
\end{equation}
Choosing $\lambda=v_i/a_i$ proves
$\Pr(X_i\leq0)\leq v_i/(v_i+a_i^2)$, including the zero-variance
case. The union bound gives Equation~\ref{eq:directional_reliability}
without requiring the $X_i$ to be independent. These random variables
usually share the same $A$ and are dependent.

The matrix $K$ is a Gram matrix in the inner product induced by $D$;
it is positive semidefinite, but $Kw$ can have nonpositive entries.
The signs of those entries determine expected task progress for the
chosen mixture. Neither an off-diagonal negative entry nor a local
nonpositive margin establishes that the capability targets are globally
unattainable. The analysis also permits successful trajectories with
temporary task regressions.

\subsection{A sufficient local sampling budget}
\label{app:budget_reliability}

Assume $a_i>0$ for every domain. For positive fractions $q_j$ summing
to one, suppose the counts $m_j=Gq_j$ are integers and satisfy any
applicable bounds. Define
\begin{equation}
    B_i(q)=\sum_j
        \frac{w_j^2\mu_i^\top D\Sigma_jD\mu_i}{q_j a_i^2}.
\end{equation}
Then $v_i(m)/a_i^2=B_i(q)/G$. For any $0<\delta<1$, the sufficient
condition
\begin{equation}
    G\geq\frac{M-\delta}{\delta}\max_i B_i(q)
    \label{eq:reliability_budget}
\end{equation}
ensures $\Pr(\exists i:X_i\leq0)\leq\delta$.
Indeed, each Cantelli term is $B_i/(G+B_i)$, which is at most
$\delta/M$ under the stated condition; summing proves the claim.
If all $B_i=0$, the sampled first-order progress is deterministic.

This is a sufficient condition on known population quantities. Rounding
counts requires recomputing the bound for the resulting allocation.
Count caps may prevent the budget from being attained. The condition
also omits estimation uncertainty, finite-population corrections,
curvature, and state evolution; it is not an operational guarantee for
AdamW or a requirement for long-run capability integration.

\subsection{Trace variance and exact integer allocation}
\label{app:allocation_proofs}

By Equation~\ref{eq:combined_covariance},
\begin{equation}
    V(m)=\operatorname{tr}(D\Sigma(m)D)
      =\sum_i\frac{w_i^2\operatorname{tr}(D\Sigma_iD)}{m_i}
      =\sum_i\frac{w_i^2e_i}{m_i}.
\end{equation}
For the positive semidefinite matrix $Q=D\Sigma(m)D$,
$\mu_i^\top Q\mu_i\leq\|\mu_i\|_2^2\lambda_{\max}(Q)
\leq\|\mu_i\|_2^2\operatorname{tr}(Q)$, proving
Equation~\ref{eq:trace_directional_bound}.
Equation~\ref{eq:noise_estimate} is the sum of coordinate-wise sample
variances and is unbiased for $e_i$ under independent sampling.
The moving average is an allocation estimate, not an unbiased estimate
of a changing current-state covariance.

Write $c_i=w_i\sqrt{e_i}$. Cauchy--Schwarz gives
\begin{equation}
    \left(\sum_i m_i\right)
    \left(\sum_i\frac{c_i^2}{m_i}\right)
      \geq\left(\sum_i c_i\right)^2.
\end{equation}
When all $c_i>0$, equality holds at
$m_i=Gc_i/(\sum_jc_j)$, yielding the continuous per-step rule without
count bounds. Positive lower and upper bounds lead to the corresponding clipped
continuous rule by the Karush--Kuhn--Tucker conditions. Zero-noise
components are handled by the bounds and remaining budget.
We do not claim that rounding this continuous solution solves the
integer problem. Both GPAS modes evaluate their estimated objectives
on every $m\in\mathcal M$ as defined in
Appendix~\ref{app:implementation}.

With positive marginal times and zero fixed step time,
\begin{equation}
    \left(\sum_i m_i\tau_i\right)
    \left(\sum_i\frac{c_i^2}{m_i}\right)
      \geq\left(\sum_i c_i\sqrt{\tau_i}\right)^2,
\end{equation}
with equality at $m_i\propto c_i/\sqrt{\tau_i}$ when this continuous
allocation is feasible. Positive fixed time and integer bounds are
handled by enumeration. As $C$ grows on the finite feasible set,
minimizing $CV(m)+(\sum_i m_i\tau_i)V(m)$ eventually selects a
minimum-variance allocation, with time breaking any remaining variance
ties. The product criterion alone does not establish a learning-rate
or time-to-capability improvement.

For the study's Uniform counts of four and cap of eight,
$V(m)\geq\sum_i w_i^2e_i/8=V_{\rm Uniform}/2$.
This proves the stated factor-of-two ceiling under the independent,
fixed-state model for any feasible allocation. All these statements
concern the same means and covariances, not different training states
reached by different schedules.

\subsection{Trace-optimal allocation need not preserve task directions best}
\label{app:trace_counterexample}

Consider two tasks with $D=I$, $w_1=w_2=1/2$, and
$\mu_1=\mu_2=(1,0)^\top$. Let
\begin{equation}
    \Sigma_1=\begin{pmatrix}1&0\\0&80\end{pmatrix},\qquad
    \Sigma_2=\begin{pmatrix}9&0\\0&0\end{pmatrix}.
\end{equation}
Both progress margins equal one. For $G=40$ with positive counts and
no additional caps, the trace noises $(e_1,e_2)=(81,9)$ give the
GPAS optimum $(m_1,m_2)=(30,10)$ and
\begin{equation}
    v_1=v_2=\frac14\left(\frac1{30}+\frac9{10}\right)
      =\frac7{30}.
\end{equation}
The allocation $(10,30)$ instead gives
$v_1=v_2=\tfrac14(1/10+9/30)=1/10$.
The latter allocation minimizes the common directional variance:
its relevant standard deviations are in the ratio $1:3$.
The trace variances are respectively $9/10$ and $21/10$, so the
trace and directional objectives prefer different allocations.

If the independent micro-batch noise is Gaussian, then
$X_1=X_2$ is Gaussian with mean one and the stated variance.
Its exact non-descent probability is $\Phi(-1/\sqrt{v_1})$, where
$\Phi$ is the standard normal distribution function. Thus $(10,30)$
also has lower exact failure probability, not merely a lower bound.
The discrepancy comes from task 1's second-coordinate variance, which
is orthogonal to both task means. This is an analytical counterexample;
its two-task budget is separate from the four-domain training protocol.

\subsection{Finite-step progress and the optimizer approximation}
\label{app:descent}

Suppose $L_i^t$ is $\mathcal L_i$-smooth along every update segment
under consideration. For the idealized update,
\begin{equation}
    L_i^t(\theta_t-\eta DA)-L_i^t(\theta_t)
      \leq-\eta X_i+
         \frac{\eta^2\mathcal L_i}{2}\|DA\|_2^2.
    \label{eq:finite_step_descent}
\end{equation}
Taking expectations gives
\begin{equation}
    \mathbb E[\Delta L_i^t]
      \leq-\eta a_i+\frac{\eta^2\mathcal L_i}{2}
          \left(\|D\mu_w\|_2^2+V(m)\right).
    \label{eq:expected_task_descent}
\end{equation}
Consequently, when all margins are positive, a sufficient condition
for every expected task-loss change to be negative is
\begin{equation}
    0<\eta<\min_i
       \frac{2a_i}{\mathcal L_i(\|D\mu_w\|_2^2+V(m))},
\end{equation}
with an infinite threshold for a zero smoothness constant.
A realized first-order descent sign by itself is insufficient:
Equation~\ref{eq:finite_step_descent} additionally requires
$X_i>\eta\mathcal L_i\|DA\|_2^2/2$ to certify a strict decrease.

A probability version requires control of update magnitude as well.
If $\Pr(\|DA\|_2>R)\leq\varepsilon_R$ and
$d_i=a_i-\eta\mathcal L_iR^2/2>0$ for all $i$, Cantelli and a union
bound give
\begin{equation}
    \Pr(\exists i:\Delta L_i^t\geq0)
       \leq\min\left\{1,\varepsilon_R+
             \sum_i\frac{v_i(m)}{v_i(m)+d_i^2}\right\}.
\end{equation}
The event of a nondecrease with $\|DA\|_2\leq R$ implies
$X_i\leq\eta\mathcal L_iR^2/2$ for some $i$, which proves the bound.
These assumptions need to be verified before interpreting the expression
as a finite-step guarantee.

A second-order expansion states how covariance direction and curvature
enter more precisely. Let $H_i=\nabla^2L_i^t(\theta_t)$. If the Hessian
is $\kappa_i$-Lipschitz on all update segments and
$\mathbb E\|DA\|_2^3<\infty$, then
\begin{equation}
\begin{split}
    \mathbb E[\Delta L_i^t]
       ={}&-\eta a_i+\frac{\eta^2}{2}
          \left[(D\mu_w)^\top H_i(D\mu_w)
          +\operatorname{tr}(H_iD\Sigma(m)D)\right]+\mathcal E_i,\\
    |\mathcal E_i|\leq{}&\frac{\kappa_i\eta^3}{6}\mathbb E\|DA\|_2^3.
\end{split}
\end{equation}
This follows by Taylor's theorem and the quadratic-form
bias--variance identity. The curvature-weighted covariance term can
have either sign when the Hessian is indefinite. The smoothness bound
replaces this directional dependence with the conservative trace term.

For the weighted frozen objective $F_t$, the same argument gives, with
$\mathcal L$ a smoothness constant,
\begin{equation}
    \mathbb E F_t(\theta_t-\eta DA)
       \leq F_t(\theta_t)-\eta\mu_w^\top D\mu_w
       +\frac{\eta^2\mathcal L}{2}
         \left(\|D\mu_w\|_2^2+V(m)\right).
\end{equation}
Improving this weighted bound does not ensure progress for every task.
The idealized update is also not the actual AdamW update: momentum,
current-gradient second moments, clipping, and decay can change its
mean and direction. Allocation can alter these states over training.
The common-checkpoint experiment therefore compares frozen-$D$
predictions with actual optimizer branches using identical saved states.

\subsection{Conditional extension: allocation for directional uncertainty}
\label{app:directional_extension}

If measurements show that trace noise is a poor proxy for relevant
uncertainty, a possible extension is to choose counts by
\begin{equation}
    \min_{m\in\mathcal M}\max_{1\leq i\leq M}
           \frac{v_i(m)}{a_i^2},\qquad a_i>0\quad\text{for every }i.
    \label{eq:directional_extension}
\end{equation}
With known positive margins and covariances, each normalized variance
is a sum of nonnegative multiples of $1/m_j$. It is convex on positive
counts; their maximum is convex, and the sum and box constraints are
convex in the continuous relaxation. With four tasks the integer
problem can again be solved by finite enumeration. This criterion
minimizes the worst normalized directional variance, and hence the
worst Cantelli bound, under the stated local model. It need not minimize
the true joint failure probability or long-run capability loss.

If any margin is nonpositive, the criterion does not resolve joint
progress. Excluding that domain would conceal the unresolved task
tradeoff. Changing weights, the update direction, or parameter sharing
requires a separate intervention. Estimating directional covariances
also adds computation: independent calibration and trial samples, or
historical estimates fixed before the trial, are needed to test the
predictions without reusing the same noise optimistically. The proposed
study first tests GPAS and its trace statistic. This extension is a
research hypothesis, not an implemented or validated algorithm.

\subsection{Timing and planned training records}
\label{app:system_details}

Task blocks record rollout, teacher-scoring, backward, and statistics
time, together with response lengths and peak memory. The marginal
per-micro-batch estimates form $\tau_i$; synchronization, optimizer
work, and other count-independent step operations form $C$.
Their EMA estimates use decay $0.9$. We compare predicted and
observed step times and evaluate efficiency with actual elapsed
time on both reserved GPUs. Evaluation compute is reported separately.
If a deployment overlaps task blocks, its timing predictor should
represent that execution schedule.

Checkpoints retain optimizer state, task stream positions, counts,
and noise and time averages. The planned traces include per-task
held-out loss, prompt consumption, response length and truncation,
micro-batch counts, noise estimates, and measured step time.
\missing{Add the artifact location and populate the configuration
and empirical figures from these records.}
